\documentclass[11pt]{exam}
%\usepackage[portuges]{babel}
\usepackage{enumerate}
\usepackage[utf8]{inputenc}
\usepackage{algorithmicx}
\usepackage{algpseudocode}

\usepackage{alltt}
\usepackage{xspace}
\usepackage{longtable}
\usepackage{amsfonts}
\usepackage{latexsym}
\usepackage{amsmath}
\usepackage{amssymb}
%\usepackage{xcolor}
\usepackage{enumerate}
\usepackage{graphicx}
%\usepackage{moreverb}
%\usepackage{ifthen}
\usepackage{url}
\usepackage{tikz} 
\usepackage{bussproofs}


\input sli_aux
\newboolean{portuguese}
\setboolean{portuguese}{true}
\setboolean{portuguese}{false}

\newboolean{normal}
\setboolean{normal}{true}

\newcommand{\pten}[2]{\ifthenelse{\boolean{portuguese}}{#1}{#2}}
\pten{\usepackage[portuges]{babel}}{\usepackage[british]{babel}}

\setcounter{aula}{5}

\begin{document}
\begin{center} 
{\bf {\large \pten{Programação Concorrente}{Concurrent Programming} - \pten{Exercícios}{Exercises} \theaula} }\\

\medskip
\pten{ Bissimilaridade Fraca e Congruência Observável (ou Equivalência Observável)  }{Weak bisimilarity and observable congruence}
\end{center}

\begin{questions}
	\question \pten{Sendo}{Let}
\begin{eqnarray*}
	Buffer &:=& put?. get?.Buffer\\
	BufferL &:=& put?. pass!.BufferL\\
	BufferR &:=& pass?. get?.BufferR
\end{eqnarray*}
\pten{Mostra que}{Show} 
$(BufferL| BufferR)\backslash\{pass!,pass?\}\approx (Buffer|Buffer)$.
\question\pten{ Sendo a bissimilaridade fraca}{Given the weak bisimilarity} $\approx$ \pten{dado por}{} 
$$\approx =\bigcup_{\textrm{R \pten{é bissimulação fraca}{weak bisimulation}}}R$$
\pten{Mostra que}{Show that}
\begin{parts}
\part \pten{A relação}{}	 $\approx$ \pten{é de equivalência}{is an equivalence}.
\part \pten{A relação}{}	 $\approx$ \pten{é maior bissimulação fraca}{is the largest weak bisimulation}.
\part \pten{A bissimilaridade fraca é estritamente mais grosseira que a bissimilaridade forte}{The weak bisimilarity is coarsest than the strong bisimilarity}, i.e.,  $\sim\subsetneq \approx$.
\end{parts}

\question
\pten{Seja}{Let} 
\begin{eqnarray*}
	CMB&:=& coin?.coffee!.CMB+coin?.CMB\\	
CS&:=&pub!.coin!.coffee?.CS\\
UniBad&:= & (CMB|CS)\backslash \{coin,cofee\}\\
\end{eqnarray*} 
\pten{Será que}{Is true that} $Spec\approx Unibad$ \pten{onde}{where} $Spec:=pub!.Spec$?
\question \pten{Mostra que para qualquer}{Show that for any} $P,Q\in CCS$ \pten{e}{and} $\alpha\in Act$ \pten{as seguintes equivalências verificam-se (As chamas leis $\tau$ de Milner)}{the following equivalences hold}:
\begin{eqnarray*}
\alpha.\tau. P &\approx& \alpha.P\\
	P+ \tau. P &\approx& \tau.P\\
\alpha.(P+\tau.Q) &\approx& \alpha.(P+\tau.Q)  + \alpha.Q\\
\end{eqnarray*}

\pten{Nota: constrói bissimulações fracas apropriadadas}{}.

\question \pten{Considera o seguinte sistema de transições}{Considere the following LTS} \pten{e}{and} \pten{mostra que}{show that} $s\not\approx t$ 
\pten{mostrando que um jogo de bissimilaridade fraca onde o atacante tem uma estratégia universal ganhadora, isto é, ganha para todas as possíveis escolhas do defesa}{using a weak bisimulation game}.

\begin{tabular}{lccr}
\begin{tikzpicture}[>=stealth, shorten >=2.5pt, auto, node distance=1.5cm, initial text={}]
\node[state,  inner sep=1pt, minimum size=5pt,  initial] (S0){$s$};
\node[state,  inner sep=1pt, minimum size=5pt][below of =S0] (S1){$s_1$};
\node[state,  inner sep=1pt, minimum size=5pt][below  of =S1] (S2){$s_2$};
\node[state,  inner sep=1pt, minimum size=5pt][below right of =S1] (S3){$s_3$};

\path[->](S0) edge  node {$a$} (S1)
	(S1) edge  node {$a$} (S2)
	(S1) edge  node {$b$} (S3);
\end{tikzpicture}&&
\begin{tikzpicture}[>=stealth, shorten >=2.5pt, auto, node distance=1.5cm, initial text={}]
\node[state,  inner sep=1pt, minimum size=5pt,  initial] (S0){$t$};
\node[state,  inner sep=1pt, minimum size=5pt][below  of =S0] (S1){$t_1$};
\node[state,  inner sep=1pt, minimum size=5pt][below left of =S1] (S2){$t_2$};
\node[state,  inner sep=1pt, minimum size=5pt][below right of =S1] (S3){$t_3$};
\node[state,  inner sep=1pt, minimum size=5pt][below of =S3] (S5){$t_5$};
\node[state,  inner sep=1pt, minimum size=5pt][below  of =S2] (S4){$t_4$};
\path[->](S0) edge  node  {$\tau$} (S1)
(S1) edge  node  {$a$} (S2)
(S1) edge  node  {$a$} (S3)
(S2) edge  node  {$a$} (S4)
(S3) edge  node  {$\tau$} (S2)
(S3) edge  node  {$b$} (S5)
;
\end{tikzpicture}
\end{tabular}

\question \pten{Considera o seguinte protocolo}{Consider the following protocol} $Protocol:= acc?.del!.Protocol$ \pten{que corresponde a um canal em que se envia e recebe mensagens (como um buffer duma posição)}{that corresponds to a comunication channel}. \pten{Seja a seguinte implementação}{Given the implemnetation}
\begin{eqnarray*}
	Send&:=& acc?.Sending\\
Sending&:=& send!. Wait\\
Wait &:= &ack?.Send + error?.Sending\\
Rec &:=& trans?.Del\\
Del &:=& del!.Ack\\
Ack &:=& ack!.Rec\\
Med &:= &send?.Med1\\
Med1 &:=& \tau.Err+trans!.Med\\
Err &:=& error!.Med
\end{eqnarray*}
\pten{Mostra que}{show that} $(Send|Med|Rec)\backslash\{send, error,trans,ack\}\approx Protocol$.
\pten{Verifica no}{Check with} Pseuco.Com.
\question \pten{Mostra que}{Show that} $a.0+0\not\approx a.0+\tau.0$ \pten{e}{and} \pten{concluí que}{conclude that} $\approx$ \pten{não é uma congruência em }{is not a congruence in} $CCS$.
\question \pten{Sendo}{Let} $\simeq$ \pten{a congruência observável, mostra que}{be the observable congruence; show that:}
\begin{itemize}
	\item $\tau.a\not\simeq a$
\item $P|\tau.Q\not\simeq \tau.(P|Q)$ 
\end{itemize}	
\question 
\pten{A seguir apresentam-se vários algoritmos }{Consider the following algorithms} $\mathit{MUTEX}$ \pten{para a exclusão mútua}{for mutual exclusivity}. \pten{Para cada um deles}{For each one}:
\begin{parts}
\part \pten{Implementa o algoritmo directamente em CCS explicando sucintamente os processos usados e qual o processo que corresponde ao algoritmo (que deverá ter o nome correspondente)}{Write it in CCS}. \pten{Assinala a zona crítica com as ações}{Use} $enter_i$ \pten{e}{and} $exits_i$ \pten{ para}{for} $i=1,2$ \pten{}{to indicate the critical region}.
	\part \pten{Fornece um ficheiro correspondente para o pseuco.com e testa para verificar (traços aleatórios) que realmente a exclusão mútua ocorre (ou não). Qual o problema caso não ocorra?}{Check if the mutual exclusion seems to be working by randomly navigating through pseuco.com.}
\part\pten{ Implementa-os também em}{Implement them also in} CCS-CAOS.
\part 
\pten{Supõe que se modela a  entrada e a saída da zona crítica por}{Consider now:}
\begin{eqnarray*}
	\mathit{MutexSpec} &:=& \mathit{enter1.exit1.MutexSpec + enter2.exit2.MutexSpec} 
\end{eqnarray*}	
\pten{Será verdade que}{Is true that}  $\mathit{MUTEX}\approx \mathit{MutexSpec}$?
%\part (***) Implementa o algoritmo na linguagem Pseuco (e/ou CAAL)  e testa também a exclusão mútua usando  traços aleatórios.

\end{parts}
\begin{enumerate}[i)]
\item \pten{Modelar o algoritmo de Peterson com dois processos para a exclusão mútua em CCS}{Peterson algorithm}.

\pten{Para o processo}{For process} $P_i$, $j,i=1,2$ \pten{e}{and} $i\not=j$.

		 \begin{algorithmic}
		\While{\bf true}
	\State {\color{green}noncricital actions}
\State{$b_i\gets \mathbf{true}$;}
\State{$k\gets j$;}
\While{$b_j \land k=j$} \State{\textbf{skip};}
\EndWhile
	\State {\color{red}critical actions}
\State$b_i\gets \mathbf{false}$
		\EndWhile		
\end{algorithmic}

\item 
\pten{Considera o algoritmo de Hyman para a exclusão mútua}{Hyman algorithm}. \pten{As variáveis}{Variables}  $b_i$ \pten{são booleanas}{are Boolean} \pten{e}{and} $k$ \pten{é inteira}{is an integer}.
\pten{Para os processos}{For processes}  $P_i$, $j,i=1,2$ \pten{e}{and} $i\not=j$.

 \begin{algorithmic}
		\While{\bf true}
	\State {\color{green!75!black}noncricital actions}
\State{$b_i\gets \mathbf{true}$;}
\While{$k\not=i$} 
\While{$b_j$}
\State{\textbf{skip};}
\EndWhile
\State{$k \gets i$;}
\EndWhile
	\State {\color{red}critical actions;}
\State{$b_i\gets \mathbf{false};$}
		\EndWhile		
\end{algorithmic}
\item 
\pten{Considera o algoritmo de Pnueli  para a exclusão mútua}{Pnueli algorithm}. 

\pten{As variáveis}{Variables}  $y_i$ \pten{são booleanas}{are Boolean} \pten{e}{and} \pten{começam em}{start in} $false$ \pten{sendo locais}{being local}. \pten{A variável}{The variable} $s$ \pten{é partilhada e toma os valores}{is shared and has value} $0$ \pten{ou}{or} $1$ \pten{começando em}{starting in} $1$.
\pten{Para o processo}{For process} $P_i$ \pten{e}{and}  $i=0,1$:
 \begin{algorithmic}
		\While{\bf true}
	\State {\color{green!75!black} noncricital actions}
\State{$y_i\gets \mathbf{true}$;}
\State{$s\gets i$;}
\While{$\lnot(y_{1-i}=0\lor (s\neq i))$}
\State{\textbf{skip};}
\EndWhile
	\State {\color{red}critical actions;}
\State{$y_i\gets \mathbf{false};$}
		\EndWhile		
\end{algorithmic}


\item 
\pten{Considera o algoritmo simplificado de exclusão mútua}{Simple algorithm}. 

\pten{A variável}{Variable}  $y$ \pten{é partilhada e booleana}{is shared and Boolean}, \pten{e}{and} \pten{começa em}{starts in} $0$.
\pten{Para o processo}{For process} $P_i$ \pten{e}{and}  $i=0,1$:
\begin{algorithmic}
\While{\bf true}
	\State {\color{green!75!black} noncricital actions}
  \While{$y=0$}
    \State{\textbf{skip};}
  \EndWhile
  % \State{$y\gets 0$; $\mathbf{crit_i}$; $y\gets 1$;}
  \State{$y\gets 0$;}
	\State {\color{red}critical actions;}
	\State{$y\gets 1$;}
\EndWhile		
\end{algorithmic}

\end{enumerate}

\question \pten{Resolve os exercícios no livro do Pseuco.com}{Solve the exercises from the Pseuco.com book}:
\begin{itemize}
	\item \url{https://book.pseuco.com/#/read/equality/weak-bisimilarity}
	\item \url{https://book.pseuco.com/#/read/equality/observation-congruence}
	\item \url{https://book.pseuco.com/#/read/equality/workout}
\end{itemize}

\end{questions}
\end{document}